\section{A wider collision band and projected window probabilities}
\label{sec:stationary-collision-band}
A prescribed collision interval does not incur a return-clock stopping
error. We exploit this distinction before passing to physical time.
It permits a full isotropic collision band with rescaled radius
$m^{9/100}$, without asserting such a band at the prescribed return
count. The initial collision density may be nonstationary, provided
that its supremum and initial-coordinate variation are bounded.
This includes the length-biased collision marginal of equilibrium flow.

Write $S_{m,R}h_R=\sum_{j=0}^{m-1}h_R\circ T_R^j$ and put
\[
 e_0=\frac3{280},\qquad
 \mathfrak e_m=m^{-e_0}\sqrt{\log(2+m)}.
\]
All initial densities in this section are relative to the collision
probability $\nu$, not to $\nu_R^*$.

\subsection{Eliminating the stopping loss on a fixed collision interval}
\begin{theorem}[An isotropic collision band with bounded-BV initial data]
\label{thm:wide-collision-band}
Fix $B,C_0<\infty$. For every complex collision function $a$ with
$\|a\|_\infty+\|a\|_{\mathrm{BV}}\le B$, and
$\alpha_a=\int a\dd\nu$, one has
\begin{align}
 &\int_{|v|\le C_0m^{9/100}}
 \left|\int a e^{iv\cdot S_{m,R}h_R/\sqrt m}\dd\nu
                   -\alpha_a e^{-v^{\mathsf T}\Gamma_Rv/2}\right|\dd v
 \le C_{B,C_0}\mathfrak e_m .
 \label{eq:wide-collision-band}
\end{align}
The bound is uniform in $R\in I$ and $m\ge2$. The physical frequencies
have radius $C_0m^{-41/100}$. There is no stopping time in the transform.
\end{theorem}
\begin{proof}
On $|v|\le2m^{1/200}$, integrate
Lemma~\ref{lem:explicit-frequency-error} with
$\delta=m^{-1/14}/4$. The six resulting terms, with
$U=m^{1/200}$ and $L_\delta=1+|\log\delta|$, are
\[
 C_B\left(\delta U^4+\sqrt{\delta L_\delta}\,U^5
 +\delta L_\delta U^6+m^{-1/2}\delta^{-4}U^5
 +m^{-1/2}\delta^{-6}U^7+\delta^{-2}\rho^mU^4\right).
\]
The slowest exponent is $1/28-5/200=3/280$. The cubic exponent
$1/2-6/14-7/200=51/1400$ and all other powers have a strict
margin over $3/280$. Thus this integral is $O_B(\mathfrak e_m)$.
No return-clock comparison is used in this calculation.

For the annulus $2m^{1/200}\le|v|\le C_0m^{9/100}$ use
Theorem~\ref{thm:high-order-damped-unsmoothing}, with a mark at the
initial collision, and the fixed choices
\begin{equation}\label{eq:wide-collision-orders}
 P=40,\quad Q=29,\quad
 \delta=m^{-19/100}/4,\quad \epsilon=m^{-50}/4.
\end{equation}
The residual Taylor degree is $79$; the spectral Taylor degree is
$29$. Neither varies with $m$. For $z=v/\sqrt m$, the two analytic
margins in \eqref{eq:finite-jet-damping-domain} are
\[
 \frac12-\frac9{100}-2\frac{19}{100}=\frac3{100},\qquad
 28\left(\frac12-\frac9{100}\right)
              -2\frac{19}{100}\,30=\frac2{25}.
\]
Fixed multiplicative constants $C_0$ are absorbed by taking $m$
sufficiently large. The volume is $O_{C_0}(m^{9/25})$.
The fine insertion and fine residual errors have decay exponents
\[
 50-\frac9{25}=\frac{1241}{25},\qquad
 50-\frac9{25}-79\left(\frac12+\frac9{100}\right)
                                      =\frac{303}{100}.
\]
For the residual moment term indexed by $1\le b\le40$, the power
before taking its negative decay sign is
\begin{equation}\label{eq:wide-collision-residual-powers}
 \frac9{25}-40+80\frac9{100}+\frac{81}{100}b.
\end{equation}
Its largest value, at $b=40$, is $-1/25$. Every smaller $b$ improves
this by $81/100$ per removed block. In particular the connected
$b=1$ term is included. The fixed logarithmic powers are absorbed
by $1/25-3/280=41/1400>0$.

The damped word has a fixed polynomial prefactor, possibly large,
and $m|z|^2\ge4m^{1/100}$. It is smaller than every fixed inverse
power after integration. The omitted Gaussian contribution has the
same property by the already proved uniform ellipticity of
$\Gamma_R=c_*D_R$. This proves the annular estimate and then
\eqref{eq:wide-collision-band}. Increasing the constant covers the
finitely many remaining $m$. The argument enlarges a collision-clock
band, not the known prescribed-return-count Fourier region.
\end{proof}

\subsection{Projection without taking a trace of a Fourier estimate}
For a positive definite matrix $C$, write $g_C$ for its normalized
Gaussian density. Put
\begin{equation}\label{eq:stationary-window-exponents}
 \sigma_1=\frac2{25},\qquad
 \kappa_1=\frac{23}{2800},\qquad
 \varepsilon_m=m^{-\kappa_1}\sqrt{\log(2+m)}.
\end{equation}

\begin{theorem}[Projected collision windows under bounded initial densities]
\label{thm:projected-collision-windows}
Let $\sigma_R\ge0$ satisfy $\int\sigma_R\dd\nu=1$ and
$\|\sigma_R\|_\infty+\|\sigma_R\|_{\mathrm{BV}}\le B$.
Let $P_{m,R}:\R^4\to\R^3$ extend to a matrix $\mathsf A_{m,R}$ with
$\|\mathsf A_{m,R}\|+\|\mathsf A_{m,R}^{-1}\|\le C_0$. Set
\[
 Y=\frac1{\sqrt m}P_{m,R}\sum_{j=0}^{m-1}h_R\circ T_R^j.
\]
For $|\xi|\le A$ and $c_0m^{-\sigma_1}\le h_i\le C_0m^{-\sigma_1}$,
\begin{align}
 \int\sigma_R\one_{\{Y\in B_3(\xi,\mathbf h)\}}\dd\nu
 =\int_{B_3(\xi,\mathbf h)}
       g_{P_{m,R}\Gamma_RP_{m,R}^{\mathsf T}}(y)\dd y
       +O_{A,B,c_0,C_0}\left(V_3\varepsilon_m\right),
 \label{eq:projected-collision-windows}
\end{align}
where $V_3=\prod_{i=1}^3(2h_i)$. Endpoints in the exact event may
be included or excluded. No smoothing is added to that event.
\end{theorem}
\begin{proof}
Write $\mathsf A=\mathsf A_{m,R}$. Use the full
four-dimensional variable $Y'=\mathsf A(\sum_{j<m}h_R\circ T_R^j)/\sqrt m$.
Truncate its fourth coordinate to $[-T,T]$, with $T=m^{1/400}$.
Choose the same bandwidth $b=c(C_0)m^{9/100}$ in all four output
coordinates. The linear change of frequency $v\mapsto\mathsf A^{\mathsf T}v$
maps this cube into the band of Theorem~\ref{thm:wide-collision-band},
with bounded determinant factors. The characteristic error integrated
on the cube is $O(\mathfrak e_m)$.

Apply the endpoint-safe product envelopes
\eqref{eq:box-lower-envelope} to
$B_3(\xi,\mathbf h)\times[-T,T]$. Their $L^1$ norms are $O(TV_3)$,
so Fourier comparison costs $CTV_3\mathfrak e_m$. To estimate the
Gaussian envelope error, use a joint density bounded by
$C\exp(-c|y|^2)$, uniformly in all parameters. In an error term for
one of the first three coordinates, integrate that error factor and
the other two output majorants in $L^1$, but integrate the fourth
majorant against the Gaussian decay. Since it is at most three,
this contributes $CV_3/(bh_i)$, without a factor $T$. The error
factor for the fourth coordinate contributes $CV_3/b$.

Lemma~\ref{lem:fixed-even-maximal} at order $128$, followed by
$\sigma_R\le B$, gives a uniform $128$th moment for $Y'_4$.
Its truncation error is at most $CT^{-128}$. The Gaussian tail
costs at most $CV_3e^{-cT^2}$. Thus the absolute error divided by
$V_3$ is bounded by
\begin{equation}\label{eq:stationary-projection-budget}
 C\left(T\mathfrak e_m+\sum_{i=1}^3(bh_i)^{-1}+b^{-1}
                  +T^{-128}/V_3+e^{-cT^2}\right).
\end{equation}
The relevant exponents are
\[
 \frac3{280}-\frac1{400}=\frac{23}{2800},\quad
 \frac9{100}-\frac2{25}=\frac1{100},\quad
 \frac{128}{400}-3\frac2{25}=\frac2{25}.
\]
They give \eqref{eq:projected-collision-windows}. Fourier comparison
has used a genuine four-dimensional integral and then removed a
coordinate by a moment bound; it has not restricted an $L^1(\R^4)$
estimate to a plane. Positivity and the endpoint envelope inequalities
apply directly even when output coordinates have atoms.
\end{proof}
